% !TeX root = ../cosmology_summary.tex

%% --------------------------------------------------------
\chapter{Ньютонівська космологія і закон збереження енергії}
\label{app:newton_energy}
%% --------------------------------------------------------

\section{Рівняння Фрідмана з ньютонівської механіки}

Один із найцікавіших фактів у космології: перше рівняння Фрідмана
можна \emph{вивести} з ньютонівської механіки~--- без жодної ЗТВ.
Звичайно, це збіг лише частковий і має глибоке пояснення
(теорема Біркгофа, про яку нижче), але педагогічна цінність
такого підходу важко переоцінити.

\begin{keybox}{Про «евристичний» статус виводу}
Слід підкреслити, що наведене нижче виведення є \emph{евристичним},
а не строгим у рамках ньютонівської теорії. Класична механіка
Ньютона, застосована до безмежного однорідного середовища,
стикається з відомими труднощами~--- \emph{парадоксом Зелігера}
(або парадоксом Неймана-Зелігера): у статичному безмежному
Всесвіті гравітаційний потенціал розбігається, а сила, що діє
на пробне тіло, стає невизначеною через умовно збіжний ряд.
Сучасне розв'язання цієї проблеми полягає в тому, що закон
всесвітнього тяжіння в космологічному контексті слід розуміти
не як дію на відстані, а як локальне рівняння поля (рівняння
Пуассона), де густина джерела $\rho$ відраховується від
середнього космологічного значення. Строге ж обґрунтування
однорідного ізотропного розширення дає лише загальна теорія
відносності, яка усуває ці парадокси шляхом динамізації
самого простору-часу. Наведений нижче вивід слід сприймати
як наочну ілюстрацію, що дозволяє вгадати форму рівнянь
Фрідмана, але не як альтернативу ЗТВ.
\end{keybox}

\textbf{Вирізаємо кулю.}
Розглянемо однорідний ізотропний Всесвіт густиною $\rho(t)$.
Вирізаємо уявну кулю радіусом $r = a(t)\,r_0$,
де $r_0$~--- фіксована супутня координата.
За теоремою Гаусса для ньютонівської гравітації,
зовнішня оболонка не діє на внутрішню.
Тому на пробну частинку масою $m$ на поверхні кулі діє лише
гравітація внутрішньої маси:
\[
  M = \frac{4}{3}\pi r^3 \rho(t).
\]

% !TeX root = ../Cosmology.tex

%% --------------------------------------------------------
%% Рисунок: ньютонівська космологія — вирізана куля
%% Потрібні пакети: tikz, xcolor (вже є в Cosmology.tex)
%% Вставка: між абзацом "Вирізаємо кулю" і "Закон збереження"
%% --------------------------------------------------------

\begin{figure}[htbp]
\centering
\begin{tikzpicture}[scale=1.1,
    force/.style={->, >=stealth, line width=1.4pt},
    dimline/.style={<->, >=stealth, gray!70, thin},
    label/.style={font=\small},
    slabel/.style={font=\scriptsize, gray!80},
]

%% --- фонові крапки (зовнішній Всесвіт) ---
\foreach \x/\y in {
    -3.2/2.1, -2.8/1.4, -2.5/2.6, -2.1/1.0, -1.8/2.8,
     3.2/2.1,  2.8/1.4,  2.5/2.6,  2.1/1.0,  1.8/2.8,
    -3.0/-1.5, -2.4/-2.2, -1.9/-1.0, -2.7/-0.7,
     3.0/-1.5,  2.4/-2.2,  1.9/-1.0,  2.7/-0.7,
    -3.5/0.3,  3.5/0.3,  -3.3/-0.4,  3.3/-0.4
}{
    \fill[gray!35] (\x,\y) circle (1.2pt);
}

%% --- зовнішня оболонка (штрихована) ---
\draw[gray!50, dashed, line width=0.6pt] (0,0) circle (3.0);

%% --- заповнення внутрішньої кулі ---
\fill[accent!8] (0,0) circle (1.9);

%% --- внутрішня куля ---
\draw[accent, line width=1.2pt] (0,0) circle (1.9);

%% --- центр мас / початок координат ---
\node[circle, fill=black, minimum size=2pt, inner sep=0] (O) at (0,0) {};
\node[label, left] at (O) {$O$};

%% --- пробна частинка на поверхні ---
\node[theoryred, circle, fill=theoryred, minimum size=3pt] (P) at (50:1.9) {};
\node[label, theoryred, above=5pt] at (P) {$m$};

%% --- вектор радіуса r ---
\draw[dimline] (O) -- node[slabel, above left, pos=0.55]
    {$r = a(t)\,r_0$} (P);

%% --- сила гравітації (до центру) ---
\draw[force, accent] (P) -- +(230:1.1)
    node[label, anchor=north west, accent] {$F_g = -\dfrac{GmM}{r^2}$};

%% --- вектор швидкості (від центру, по радіусу) ---
\draw[force, theoryred] (P) -- +(50:0.9)
    node[label, anchor=south west, theoryred] {$\dot{r} = H r$};

%% --- маса всередині ---
\node[label, accent!80, align=center] at (-0.3, -1)
    {$M = \dfrac{4}{3}\pi r^3\rho$};

%% --- позначення зовнішньої оболонки ---
\draw[dimline] (O) -- node[slabel, below right, pos=0.7]
    {$R_\text{ext}$} (310:3.0);

%% --- підпис "зовнішня оболонка не діє" ---
\node[slabel, align=center, text width=2.8cm] at (3, 3)
    {зовнішня оболонка \\ не діє (теор.~Гауса)};
\draw[->, gray!50, thin] (2.1, 2.15) -- (2.55, 1.5);

%% --- позначення радіуса кулі ---
\draw[dimline] (O) -- node[slabel, below, pos=0.6]
    {$r$} (230:1.9);

%% --- підпис однорідна речовина ---
\node[slabel, accent!60] at (-1.1, 1.1) {$\rho(t)$};

\end{tikzpicture}
\caption{Ньютонівська модель однорідного Всесвіту.
Вирізаємо уявну кулю радіусом $r = a(t)r_0$
з однорідного середовища густиною $\rho(t)$.
Пробна частинка $m$ на поверхні відчуває лише
гравітацію внутрішньої маси $M$ (теорема Гауса):
зовнішня оболонка не дає внеску.
Закон збереження енергії $T + U = E$ для цієї
частинки безпосередньо відтворює перше рівняння
Фрідмана~\eqref{eq:newton_friedmann}.}
\label{fig:newton_sphere}
\end{figure}


\textbf{Закон збереження енергії.}
Повна механічна енергія частинки на поверхні:
\[
  \frac{1}{2}m\dot{r}^2 - \frac{GmM}{r} = E = \text{const}.
\]
Підставляємо $r = a r_0$, $\dot{r} = \dot{a} r_0$,
$M = \frac{4}{3}\pi a^3 r_0^3 \rho$:
\[
  \frac{1}{2}m\dot{a}^2 r_0^2 - \frac{4\pi G}{3}\,m\rho\,a^2 r_0^2 = E.
\]
Ділимо на $\frac{1}{2}m a^2 r_0^2$ і вводимо $H = \dot{a}/a$:
\begin{equation}\label{eq:newton_friedmann}
  \boxed{H^2 = \frac{8\pi G}{3}\,\rho - \frac{2E}{m r_0^2 a^2}.}
\end{equation}
Це і є \textbf{перше рівняння Фрідмана}~\eqref{eq:F1}
з ідентифікацією $k = -2E/(m r_0^2)$\footnote{%
На перший погляд $k$ залежить від радіуса довільної
кулі $r_0$~--- але це не так. Повна енергія пробної
частинки $E \propto m r_0^2$ (кінетична та потенціальна
енергії обидві пропорційні $r_0^2$), тому відношення
$E/(m r_0^2)$ не залежить від вибору $r_0$. Параметр
$k$ є глобальною характеристикою геометрії Всесвіту,
а не залежить від того, яку кулю ми вирізали.
У ньютонівській картині $k$ --- це просто знак повної
енергії на одиницю маси. У ЗТВ той самий параметр
набуває геометричного змісту: він визначає знак
просторової кривини і приймає лише значення $-1, 0, +1$.
}:
\begin{itemize}
  \item $E < 0$ (зв'язана система) $\leftrightarrow$ $k = +1$ (замкнений Всесвіт),
  \item $E = 0$ (параболічна орбіта) $\leftrightarrow$ $k = 0$ (плаский Всесвіт),
  \item $E > 0$ (незв'язана система) $\leftrightarrow$ $k = -1$ (відкритий Всесвіт).
\end{itemize}

\textbf{Друге рівняння Фрідмана.}
Другий закон Ньютона для тієї ж частинки:
\[
  m\ddot{r} = -\frac{GmM}{r^2} = -\frac{4\pi G}{3}\,m\rho\,r.
\]
Підставляємо $r = ar_0$:
\begin{equation}\label{eq:newton_acceleration}
  \frac{\ddot{a}}{a} = -\frac{4\pi G}{3}\,\rho.
\end{equation}
Порівняємо з релятивістським~\eqref{eq:F2}:
\[
  \frac{\ddot{a}}{a} = -\frac{4\pi G}{3}\,(\rho + 3p) + \frac{\Lambda}{3}.
\]
Різниця принципова: у ньютонівській версії \emph{немає тиску}.
У ЗТВ тиск гравітує нарівні з масою~--- саме тому
в сильних полях і при великих швидкостях ньютонівська
космологія хибує. Для нерелятивістської матерії
($p \ll \rho c^2$) обидва виразу збігаються.

\begin{keybox}{Небезпечний збіг для фотонів}
Якщо продиференціювати рівняння~\eqref{eq:newton_friedmann}
по часу і підставити $\rho \propto a^{-4}$ (радіація),
формально отримаємо рівняння прискорення
$\ddot{a}/a = -\frac{8\pi G}{3}\rho$
(пор.\ з~\eqref{eq:newton_acceleration})~---
що точно збігається з релятивістським результатом для
випромінювання $-\frac{4\pi G}{3}(\rho + 3p) =
-\frac{8\pi G}{3}\rho$ при $p = \rho/3$.

Це \emph{математичний збіг}, а не фізична закономірність.
У ньютонівській механіці однорідний тиск не має
гравітаційної маси~--- $\nabla p = 0$ у однорідному
Всесвіті, тому тиск взагалі не впливає на прискорення
кулі. Член $3p$ у релятивістському рівнянні~--- суто
ефект ЗТВ: тиск безпосередньо викривляє простір-час
через компоненту тензора енергії-імпульсу $T^{ii}$.
Збіг числових коефіцієнтів є випадковим наслідком
диференціювання закону збереження~--- не більше.

\textit{Висновок:} ньютонівський підхід є строго
легітимним лише для пилу ($p = 0$). Для радіації та
інфлатона він дає правильну відповідь лише «за
збігом обставин».
\end{keybox}

\textbf{Чому це працює? Теорема Біркгофа.}
У ЗТВ існує точний аналог теореми Гаусса~---
\emph{теорема Біркгофа}: метрика поза сферично-симетричним
розподілом маси є метрикою Шварцшильда незалежно від того,
чи пульсує ця маса. Це гарантує, що зовнішня оболонка справді
не впливає на внутрішню, і ньютонівський аргумент виявляється
точним у ЗТВ --- аж до ролі тиску і космологічної сталої.

%% --------------------------------------------------------
\section{Закон збереження енергії в космології}
%% --------------------------------------------------------

\subsection*{У ньютонівській картині}

У виводі вище ми записали $E = \text{const}$ --- і це є
буквальним законом збереження повної механічної енергії.
Рівняння~\eqref{eq:newton_friedmann} є нічим іншим як
$T + U = E$.

Однак при розширенні Всесвіту виникає парадокс: звідки береться
кінетична енергія розльоту галактик? Відповідь ньютонівська:
за рахунок гравітаційної потенціальної енергії.
При $k = 0$ повна енергія нульова: кінетична і потенціальна
точно компенсують одна одну.

\subsection*{У ЗТВ: рівняння збереження і його статус}

Третє рівняння Фрідмана:
\[
  \dot\rho = -3H(\rho + p)
\]
виглядає як рівняння неперервності, але \emph{не є} законом
збереження енергії у звичайному сенсі. Щоб це зрозуміти,
запишемо його через об'єм $V \propto a^3$:
\[
  \dot\rho + 3H\rho = -3Hp \quad\Rightarrow\quad
  \frac{d(\rho V)}{dt} = -p\,\frac{dV}{dt}.
\]
Ліва частина --- зміна повної енергії $E = \rho V$,
права --- робота тиску $-p\,dV$.
Це виглядає як перший закон термодинаміки $dE = -p\,dV$,
і в цьому сенсі енергія «зберігається». Однак ця аналогія
є математичною тотожністю~--- не причинно-наслідковим
зв'язком. По-перше, права частина $-p\,dV$ не означає що
«простір виконує роботу над речовиною»: це лише геометричний
наслідок розширення метрики. По-друге, у космології немає
«зовнішнього середовища»~--- Всесвіт не розширюється
\emph{у щось}, тому класична термодинамічна інтерпретація
роботи тиску тут принципово неповна. Фізична інтерпретація
залежить від природи компоненти~--- як ми побачимо нижче.

Але є тонкість: \textbf{для фотонів} ($p = \rho/3$,
$\rho \propto a^{-4}$) повна енергія
$E = \rho V \propto a^{-4} \cdot a^3 = a^{-1}$
\emph{спадає}. Куди йде ця енергія? Права частина
$-p\,dV < 0$: фотони виконують додатну роботу над
розширюваним об'ємом~--- і ця робота \emph{нікуди
не акумулюється}. Вона просто розчиняється у геометрії
простору, що розширюється.

Протилежна ситуація~--- \textbf{темна енергія}
($p = -\rho$, $w = -1$): права частина
$-p\,dV = +\rho\,dV > 0$, тобто при розширенні
повна енергія $E = \rho V \propto a^3$ \emph{зростає}.
Від'ємний тиск означає, що при розширенні об'єму
$dV > 0$ права частина $-p\,dV = +\rho\,dV > 0$~---
повна енергія зростає. Але це не «робота над зовнішнім
середовищем»: у космології немає зовнішнього середовища.
Математично це просто відображення того факту що
густина $\rho \approx \text{const}$ при зростаючому
об'ємі. Саме це зростання повної енергії вакууму
і є двигуном прискореного розширення.

Для \textbf{інфлатона} під час інфляції ($w \approx -1$)
діє та сама логіка що і для темної енергії: права
частина $-p\,dV \approx +\rho\,dV > 0$, тому повна
енергія $E = \rho V \propto a^3$ \emph{зростає} разом
з об'ємом. Але тут є принципова відмінність від
фотонів, яка знімає уявний парадокс: енергія інфлатона
росте не тому, що простір «передає» йому роботу
(гравітація не виконує роботу над полем у класичному
розумінні). Вона росте тому, що густина потенціальної
енергії поля $V(\phi) \approx \text{const}$ є
фундаментальною властивістю самого вакуумного стану~---
при розширенні створюється новий простір, який
автоматично заповнюється тим самим полем з тією самою
густиною енергії. Тому «накопичена» енергія плазми
після розігріву походить не з роботи розширення, а
з потенціалу скалярного поля~--- він і є справжнім
джерелом енергії. Після закінчення інфляції інфлатон
розпадається на звичайні частинки (розігрів), і саме
ця потенційна енергія поля переходить у гарячу плазму.

\subsection*{Теорема Нетер і відсутність симетрії}

У класичній механіці закон збереження енергії є наслідком
\textbf{теореми Нетер}: симетрії лагранжіана відносно
часових зсувів ($t \to t + \varepsilon$) відповідає
збережена величина~--- енергія.

У розширюваному Всесвіті ця симетрія \emph{порушена}:
метрика FLRW явно залежить від часу через $a(t)$.
Тому глобального закону збереження енергії немає~---
і це не математична вада формалізму, а фізична реальність.

\begin{keybox}{Чи порушується закон збереження енергії у космології?}
Точніше питання: \emph{чи існує глобальна енергія яка б
зберігалась?} Відповідь~--- ні, і не тому що енергія
«порушується», а тому що у викривленому просторі-часі
FLRW \emph{немає способу однозначно визначити} таку
величину. У загальній відносності закон збереження~---
це локальне рівняння $\nabla_\mu T^{\mu\nu} = 0$,
яке завжди виконується і дає рівняння~\eqref{eq:conservation_law}.
Глобальна ж енергія не визначена бо немає єдиного
вектора Кілінга часової трансляції\footnote{%
\emph{Вектор Кілінга}~--- це векторне поле $\xi^\mu$,
що задовольняє рівнянню $\nabla_{(\mu}\xi_{\nu)} = 0$.
Фізичний зміст: поток вздовж $\xi^\mu$ є симетрією
метрики, а збережена вздовж геодезичної величина
$p_\mu \xi^\mu = \text{const}$ є відповідним інтегралом
руху. Для статичного простору-часу існує вектор Кілінга
$\xi^\mu = (1,\mathbf{0})$, якому відповідає збережена
енергія. Метрика FLRW залежить від часу через $a(t)$~---
тому такого вектора немає, і глобальна енергія не
визначена.} для метрики FLRW
(він існує лише для статичних просторів-часів).

Отже, фотони справді «червоніють» і втрачають енергію.
Ця енергія не переходить нікуди~--- вона просто зникає
у геометрії простору, що розширюється.
Саме тому космологія вимагає відмовитися від деяких
інтуїтивних уявлень класичної фізики.
\end{keybox}

\subsection*{Практичний підсумок}

\begin{itemize}
    \item \textbf{Ньютонівське виведення рівняння Фрідмана}
          є коректним лише за умови постулювання
          $r(t) \propto a(t)$ і дає правильну структуру
          лише для пилу ($p=0$).
    \item \textbf{Параметр кривини $k$} виникає як
          $-2E/(mr_0^2)$ і не залежить від вибору $r_0$;
          у ЗТВ він набуває лише значень $-1, 0, +1$.
    \item \textbf{Друге рівняння Фрідмана в ньютонівській
          версії} не містить тиску. Збіг для радіації
          після диференціювання закону збереження є
          \emph{випадковим}, а не фізичним.
    \item \textbf{Рівняння неперервності}
          $\dot\rho = -3H(\rho+p)$ є локальним законом
          збереження $\nabla_\mu T^{\mu\nu}=0$, але не
          виражає глобального збереження енергії через
          відсутність вектора Кілінга в метриці FLRW.
    \item \textbf{Аналогія $dE = -p\,dV$}~--- це
          математичний наслідок, а не фізичний перший
          закон термодинаміки; у космології немає
          «зовнішнього середовища».
    \item \textbf{Для фотонів} ($w=1/3$) повна енергія
          $E = \rho V \propto a^{-1}$ спадає; ця енергія
          не переходить нікуди, а «розчиняється» в
          геометрії простору-часу.
    \item \textbf{Для темної енергії та інфлатона}
          ($w\approx -1$) густина енергії залишається
          сталою, тому $E = \rho V \propto a^3$ зростає.
          Джерело~--- потенціальна енергія поля $V(\phi)$,
          а не робота розширення.
    \item \textbf{Глобальна енергія у викривленому
          просторі-часі} не є однозначно визначеною
          величиною; тому твердження «енергія не
          зберігається» слід розуміти як «немає способу
          визначити глобальну енергію яка б зберігалась».
    \item \textbf{Космологічне червоне зміщення}~--- це
          ефект порівняння вимірювань у різних
          локально-інерційних системах відліку, а не
          поступова втрата енергії фотоном під час
          подорожі.
\end{itemize}
